% Ad Familiares 13.18 (ad-familiares-13-18)
% Source: https://raw.githubusercontent.com/PerseusDL/canonical-latinLit/master/data/phi0474/phi056/phi0474.phi056.perseus-lat1.xml
% Fetched by scripts/fetch_latin.py

% Dateline (Perseus): Scr. Romae, ut videtur, a. 708 (46) .

\ciceroLetterOpener{CICERO SERVIO. S.}

\ciceroSection{1} non concedam ut Attico nostro, quem elatum laetitia vidi, iucundiores tuae suavissime ad eum et humanissime scriptae litterae fuerint quam mihi. nam etsi utrique nostrum prope aeque gratae erant, tamen ego admirabar magis te, qui, si rogatus aut certe admonitus liberaliter Attico respondisses †(quod tamen dubium nobis quin ita futurum fuerit non erat), ultro ad eum scripsisse eique nec opinanti voluntatem tuam tantam per litteras detulisse. de quo non modo rogare te ut eo studiosius mea quoque causa facias non debeo (nihil enim cumulatius fieri potest quam polliceris), sed ne gratias quidem agere, quod tu et ipsius causa et tua sponte feceris ;

\ciceroSection{2} illud tamen dicam, mihi id, quod fecisti, esse gratissimum. tale enim tuum iudicium de homine eo, quem ego unice diligo, non potest mihi non summe esse iucundum ; quod cum ita sit, esse gratum necesse est. sed tamen, quoniam mihi pro coniunctione nostra vel peccare apud te in scribendo licet, utrumque eorum, quae negavi mihi facienda esse, faciam ; nam et ad id, quod Attici causa te ostendisti esse facturum, tantum velim addas, quantum ex nostro amore accessionis fieri potest, et, quod modo verebar, tibi gratias agere, nunc plane ago teque ita existimare volo, quibuscumque officiis in Epiroticis reliquisque rebus Atticum obstrinxeris, iisdem me tibi obligatum fore.
